\section{Norm Learning Without a Substrate: The Installation Problem}\label{sec:cooperative-ai}

%% ~1500 words / ~2 pages. No EFE formalism. No §3 apparatus.
%% Anchor keys: dafoe2020CooperativeAI, critchKrueger2020ARCHES,
%%   koster2022SillyRules, oldenburgZhiXuan2024BayesianRule,
%%   kulveit2025GradualDisempowerment, albarracin2025SocialPower.
%% Light forward refs: albarracin2022EpistemicCommunities, henrich2009CREDs.

Cooperative intelligence --- the capacity of agents to find mutually
beneficial joint courses of action --- is the organising concept behind a
rapidly consolidating research field. \citet{dafoe2020CooperativeAI} lay out
the agenda: cooperation rests on understanding, communication, commitment, and
norms. Of these four pillars, norms carry the heaviest load. They stabilise
expectations across agents who cannot fully observe one another's preferences,
sustain coordination under incomplete contracting, and provide the decentralised
scaffolding that makes institutions tractable. Whether the agents in question are
AI--AI, AI--human, or human--human-via-AI, a world with more capable AI systems
is a world in which norms do more work, not less~\citep{critchKrueger2020ARCHES}.
Getting their computational basis right is therefore not a peripheral concern for
cooperative AI: it is foundational.

\subsection*{Norms-as-Equilibria: What Works and What Is Missing}

The dominant treatment of norms in cooperative AI follows the game-theoretic
tradition inaugurated by Lewis and formalised by Schelling: a norm is a
self-enforcing equilibrium of behaviour sustained by mutual expectations. In this
frame, a social norm \textit{exists} when compliance and sanctioning behaviour
jointly persist at population scale under coordination or social-dilemma
pressure. The practical advantage of this framing is substantial. It interfaces
naturally with multi-agent reinforcement learning (MARL), makes the stability
conditions of norms tractable, and generates testable predictions about which
equilibria persist under evolutionary pressure. The Leibo lineage of cooperative
MARL papers, the decentralised sanctioning work of \citet{vinitsky2023sanctions},
and the Bayesian rule-induction approach of \citet{oldenburgZhiXuan2024BayesianRule}
all inherit this framing, and all deliver real scientific progress within it.

What the equilibrium framing does not deliver is a \textit{cognitive substrate}
--- a mechanistic story about what a norm is in the head of a single agent.
The equilibrium frame treats the agent's internal representation as a black box:
whatever it contains, if compliance and sanctioning persist, the norm is present.
Three observations mark the limits of this silence. First, violations of norms
\textit{feel surprising} in a way that violations of mere regularities do not;
they elicit moral outrage, not just strategic recalculation. Second, some norms
are \textit{sticky}: agents continue to follow them even when enforcement is
relaxed and even when they cannot articulate a welfare rationale. Third, and
most relevant here, agents can shift from \textit{imitating} a norm --- following
it to avoid sanctions --- to \textit{embodying} it, at which point the norm shapes
what they attend to, what they notice as salient, and what they find imaginable
as a course of action. The equilibrium account has no vocabulary for this
phenomenology, because it operates at the level of policy outputs, not generative
model structure.

The \citet{koster2022SillyRules} result makes this gap vivid in a particularly
elegant way. Adding arbitrary ``silly rules'' --- taboos that carry no welfare
justification --- into a MARL foraging environment \textit{improves} agents'
acquisition of compliance and enforcement with welfare-relevant taboos. This
result is surprising under the equilibrium frame, because silly rules add no
information about the welfare-relevant equilibrium; under a precision-substrate
account, it is exactly what one would expect. Silly rules supply additional
Bayesian evidence that a \textit{normative regime is in operation} ---
they pump up the precision of the agents' prior over the policy distribution
``follow taboos'' without specifying any particular content. Once that
precision is elevated, the agents learn welfare-relevant norms faster. The
mechanism Koster et al.\ identify is not primarily about content; it is about
the legibility of the normative substrate as such. This is the empirical
signature of a precision-shaping story, written in MARL language.

\citet{oldenburgZhiXuan2024BayesianRule} is the most cognitively rich current
treatment and the most direct dialogue partner for this paper. Their agents
Bayesianly induce shared norms as discrete obligative or prohibitive rules from
observed compliance and sanction data; they show that a \textit{shared
assumption} that a norm system is operating is sufficient to bootstrap common
knowledge of specific norms, enabling rapid convergence and stable transmission
to newcomers. This is a genuine advance: unlike equilibrium-only accounts, it
gives norms propositional structure in individual agents. But the ``shared
assumption'' itself remains unmodelled prior content. It is posited as an
external parameter of the inference problem rather than as something agents
acquire from a community, transmit with fidelity, and maintain through
coordinated attention. Representing norms as discrete inferred rules also
forecloses the graded, phenomenologically thick account of norm-internationalisation
that the sticky-norm and violation-surprise phenomena require. The rules are
the outputs of inference; the question of what \textit{installs the prior} that
makes the shared assumption credible is left open.

\subsection*{The Installation Question}

Where do an agent's priors over what-counts-as-a-norm come from? This is the
question the equilibrium and Bayesian-rule accounts both treat as exogenous.
Norms are handed to agents, or emerge from iterated play, but in neither case
does the account say how one generation's normative landscape gets written into
the next generation's starting priors. In human cognition and cultural evolution,
this is precisely the question that \textit{is} answered: credibility-enhancing
displays, prestige-biased transmission, and community-level attention allocation
are the mechanisms by which prior preferences over outcomes are shaped and
reproduced~\citep{henrich2009CREDs}. Any formal account of norm learning that
cannot speak to these mechanisms is incomplete at the level of mechanism, even
if it is complete at the level of equilibrium description.

The stakes of this incompleteness become acute when AI systems participate in
human normative life. \citet{kulveit2025GradualDisempowerment} argue that
incremental AI integration can gradually disempower human populations not through
any adversarial takeover but through structural drift: as AI systems become
load-bearing participants in economic, cultural, and institutional life, the
implicit alignment between human interests and system behaviour that comes from
human participation being generative is progressively eroded. Under the equilibrium
framing of norms, AI systems that sit at the behavioural equilibrium look
unproblematic. But if norms supervene on shared cognitive architecture --- if
their stability is grounded in the precision-shaping dynamics through which a
community installs prior preferences in its members --- then AI systems that
participate only at the policy level, without contributing to the substrate that
makes norms credible, can hollow out the normative structure even while the
behaviour looks intact. The installation question is therefore not merely a
technical gap in formal models; it is the micro-level mechanism through which
gradual disempowerment operates.

\begin{figure}[t]
\centering
\begin{tikzpicture}[
  node distance = 2.4cm and 3.2cm,
  box/.style = {
    draw, rectangle, rounded corners=3pt,
    minimum width=2.4cm, minimum height=0.9cm,
    font=\small, align=center
  },
  missing/.style = {
    draw=gray!60, rectangle, rounded corners=3pt,
    minimum width=2.4cm, minimum height=0.9cm,
    font=\small, align=center, dashed, fill=gray!10
  },
  arr/.style  = {-Stealth, thick},
  marr/.style = {-Stealth, thick, dashed, gray!70}
]
  %% Main horizontal chain
  \node[box] (world)   {World};
  \node[box, right=of world]  (policy)  {Policy $\policy$};
  \node[box, right=of policy] (action)  {Action};

  %% Off-axis community node
  \node[missing, above=1.5cm of policy] (community) {Community};

  %% Horizontal arrows
  \draw[arr]  (world)  -- (policy);
  \draw[arr]  (policy) -- (action);

  %% Missing arrow: community -> prior layer (annotated)
  \draw[marr] (community) -- node[right, font=\footnotesize, text=gray!80]
              {missing / proposed} (policy);

  %% Label under the policy node to indicate the prior layer
  \node[font=\footnotesize, below=0.15cm of policy, text=gray!60]
        {prior $\C$};

\end{tikzpicture}
\caption{The missing arrow: existing cooperative-AI norm models operate on the
horizontal chain (world $\to$ policy $\to$ action) and leave the community-to-prior
channel exogenous (dashed). This paper's contribution is to formalise the shaded
arrow: how community-level precision dynamics install an agent's prior preferences
over outcomes ($\C$), constituting the cognitive substrate of norms.}
\label{fig:missing-arrow}
\end{figure}

\subsection*{Differentiation from Albarracin et al.\ (2025)}

\citet{albarracin2025SocialPower} formalise social power as the differential
capacity to set the policy-precision parameter ($\prc$) in others' generative
models. In their account, a socially powerful agent is one who can raise or
lower how confidently another agent pursues a given policy --- shaping the
effective temperature of the policy distribution without necessarily altering
what the agent considers valuable in the first place. This is a significant
contribution: it gives power a formal active-inference definition, it engages
Bourdieu's concepts of habitus, field, and capital, and it shows how
Foucauldian disciplinary mechanisms can be read as precision-modulation
operations. Our move is one layer up, on the \textit{content} of prior
preferences ($\C$) --- that is, what an agent comes to consider valuable
in the first place, rather than how confidently it pursues a given policy.
Both layers are causally involved in the cultural reproduction of norms:
a community that controls the policy-precision parameter $\prc$ in its
members (the Albarracin layer) can enforce compliance, but a community that
controls the content of $\C$ (our layer) shapes what counts as desirable
before any enforcement becomes necessary. The prior-preference content layer
is, we argue, the substrate that cooperative-AI norm theory currently lacks.
Albarracin et al.\ provide the complementary half of the story; this paper
provides the other half, and together they constitute a two-layer account of
normativity in active-inference agents. The dynamics by which epistemic
communities stabilise or fragment around shared $\C$-vectors are developed
in~\citet{albarracin2022EpistemicCommunities}, which we build on directly
in Section~\ref{sec:formal}.

\subsection*{Roadmap}

The remainder of the paper develops the missing arrow in three steps.
Section~\ref{sec:cultural-evolution} establishes the cultural-evolution backdrop,
reviewing the mechanisms --- prestige-biased transmission, credibility-enhancing
displays~\citep{henrich2009CREDs}, and community-level attention allocation ---
through which prior preferences are differentially shaped and reproduced across
agents.
Section~\ref{sec:formal} develops the formal account: we define norms as
stabilised regions of high-precision shared $\C$, characterise the installation
channel as a community-level precision-shaping process, and derive stability
conditions for the resulting attractor landscape.
Section~\ref{sec:simulation} proposes a minimal simulation instantiating the
account in a pymdp multi-agent environment, designed to exhibit the bifurcation
between norm-convergence and echo-chamber fragmentation as a function of
the precision dynamics.
